% GENERATED FILE. Edit canonical lessons, then run npm run generate:books.
% canonical-source-hash: fnv1a64:a8f6df994cf5b2a2
% canonical-lessons: ZH-C170-guo, ZH-C170-chabei, ZH-W170-bei-cup, ZH-C170-shaozi, ZH-C170-chazi, ZH-C170-weibolu

\chapter{Pots and Pans}
\label{ch:zh-pots-and-pans}
\hlchaptermodality{170}

\begin{chapteropening}
\textbf{By the end of this chapter:} \emph{I can name a wok, a teacup, a spoon, a fork and a microwave.}

Complete the last lesson of chapter 170: I can name a wok, a teacup, a spoon, a fork and a microwave.
\end{chapteropening}

\section[guō]{guō — pot; pan; wok}

\label{lesson:ZH-C170-guo}



\begin{quote}
Before the new one: say the Mandarin for to stay up late, then the Mandarin for late at night.
\end{quote}

\subsection*{You'll want to know: guō}
\emph{guō} — \textquotedblleft{}pot; pan; wok\textquotedblright{}.

\begin{quote}
\emph{zhège guō hěn zhòng} — This wok is heavy.
\end{quote}

Its characters come later; for now it is a word for the ear.

\begin{cousinweb}
\emph{guō} is inherited from Middle Chinese \emph{kwa} (Guangyun, in Baxter's transcription), glossed 'a vessel for warming'.
\end{cousinweb}

\begin{grammarlens}[title={What you've built}]
One more word for everyday things.
\end{grammarlens}

\subsection*{Your turn}
\begin{itemize}
\raggedright
  \item \emph{Say it:} \emph{guō}
  \item \emph{Say it:} \emph{guō}, once more
  \item \emph{Say it:} \emph{zhège guō hěn zhòng} — This wok is heavy.
  \item \emph{Recall:} say the Mandarin for to stay up late, then the Mandarin for late at night, then say \emph{guō} again
\end{itemize}

\begin{tcolorbox}[breakable,colback=teal!4,colframe=teal!35!black,title={Before you move on}]
What does guō mean? (\textquotedblleft{}Pot; pan; wok\textquotedblright{}.) Say it once more.
\end{tcolorbox}
\section[\texorpdfstring{\zh{茶杯} (chábēi)}{茶杯 (chábēi)}]{\zh{茶杯} (chábēi) — teacup}

\label{lesson:ZH-C170-chabei}



\begin{quote}
Before the new one: say the Mandarin for late at night, then the Mandarin for a cooking pot.
\end{quote}

\subsection*{You'll want to know: \zh{茶杯}}
\textbf{\zh{茶杯}} — \emph{chábēi} — \textquotedblleft{}teacup\textquotedblright{}.

\begin{quote}
\textbf{\zh{我觉得这个茶杯很好看}} \emph{wǒ juéde zhège chábēi hěn hǎokàn} — I think this teacup is lovely.
\end{quote}

The new character, in pieces: \zh{杯} — eight strokes in two pieces: the tree piece on the left of \zh{机}, four strokes, and \zh{不} on the right. The next lesson puts it in your hand.

\begin{cousinweb}
\textbf{\zh{茶杯}} (\emph{chábēi}) is built from \zh{茶} 'tea' and \zh{杯} 'cup', both inherited from Middle Chinese, \emph{drae} and \emph{pwoj} (Guangyun, in Baxter's transcription).
\end{cousinweb}

\begin{grammarlens}[title={What you've built}]
One more word for everyday things.
\end{grammarlens}

\subsection*{Your turn}
\begin{itemize}
\raggedright
  \item \emph{Say it:} \emph{chábēi}
  \item \emph{Say it:} \emph{chábēi}, once more
  \item \emph{Say it:} \emph{wǒ juéde zhège chábēi hěn hǎokàn} — I think this teacup is lovely.
  \item \emph{Recall:} say the Mandarin for late at night, then the Mandarin for a cooking pot, then say \emph{chábēi} again
\end{itemize}

\begin{tcolorbox}[breakable,colback=teal!4,colframe=teal!35!black,title={Before you move on}]
What does \zh{茶杯} mean? (\textquotedblleft{}Teacup\textquotedblright{}.) Say it once more.
\end{tcolorbox}
\section[\texorpdfstring{\zh{杯} (bēi)}{杯 (bēi)}]{\zh{杯} — a tree piece, then \zh{不}}

\label{lesson:ZH-W170-bei-cup}



\begin{quote}
The word first, then the character inside it.

\begin{itemize}
\raggedright
  \item \emph{Recall:} say \emph{chábēi} — \textbf{R1}, one lesson back, before it decays
\end{itemize}

\begin{quote}
\zh{不} \zh{机}
\end{quote}

\emph{Read it:} each one aloud, with nothing said first — \emph{bù}, \emph{jī}
\end{quote}

\subsection*{Script — a tree piece, then \zh{不}}
\hlblockfigure{\detokenize{figures/ZH-W170-bei-cup-filmstrip.pdf}}{How \zh{杯} is written, stroke by stroke}

\begin{quote}
\zh{杯}
\end{quote}

\textbf{Two pieces, left before right.} Count the pieces, not the strokes.

\textbf{The left piece, four strokes: the tree piece, as on the left of \zh{机}.} A horizontal. A vertical down through it. A stroke falling to the lower left from the crossing. Then a dot to the lower right, where \zh{本} has a long stroke.

\textbf{The right piece, \zh{不}, four strokes,} as you write it: a horizontal, a stroke falling to the left, a vertical, and a dot.

4 + 4 = eight.

\subsection*{Your turn}
\begin{itemize}
\raggedright
  \item \emph{Write it:} \zh{杯}, all eight, with the model in front of you
  \item \emph{Cover:} write it again from memory, and count the pieces, not the strokes
\end{itemize}

\begin{tcolorbox}[breakable,colback=teal!4,colframe=teal!35!black,title={Before you move on}]
How many strokes in \zh{杯}? (\textbf{Eight}.) What is the right-hand piece? (\textbf{\zh{不}}.) How does the tree piece end? (\textbf{With a dot, not a long stroke}.)

Source: \href{https://raw.githubusercontent.com/chanind/hanzi-writer-data/68d10a4b21150cae5e1ebbd223eed289cf32d90c/data/\%E6\%9D\%AF.json}{Hanzi Writer Data \zh{杯}, pinned stroke paths and medians}.
\end{tcolorbox}
\section[sháozi]{sháozi — spoon; ladle}

\label{lesson:ZH-C170-shaozi}



\begin{quote}
Before the new one: say the Mandarin for a cooking pot, then the Mandarin for a teacup.
\end{quote}

\subsection*{You'll want to know: sháozi}
\emph{sháozi} — \textquotedblleft{}spoon; ladle\textquotedblright{}.

\begin{quote}
\emph{wǒ yào yí ge sháozi} — Could I have a spoon?
\end{quote}

Its characters come later; for now it is a word for the ear.

\begin{cousinweb}
\emph{sháozi} is built from \emph{sháo} 'ladle' and the noun ending \zh{子}, both inherited from Middle Chinese, \emph{dzyak} and \emph{tsiX} (Guangyun, in Baxter's transcription), which glosses the first as a drinking vessel. In the word the second loses its tone.
\end{cousinweb}

\begin{grammarlens}[title={What you've built}]
One more word for everyday things.
\end{grammarlens}

\subsection*{Your turn}
\begin{itemize}
\raggedright
  \item \emph{Say it:} \emph{sháozi}
  \item \emph{Say it:} \emph{sháozi}, once more
  \item \emph{Say it:} \emph{wǒ yào yí ge sháozi} — Could I have a spoon?
  \item \emph{Recall:} say the Mandarin for a cooking pot, then the Mandarin for a teacup, then say \emph{sháozi} again
\end{itemize}

\begin{tcolorbox}[breakable,colback=teal!4,colframe=teal!35!black,title={Before you move on}]
What does sháozi mean? (\textquotedblleft{}Spoon; ladle\textquotedblright{}.) Say it once more.
\end{tcolorbox}
\section[chāzi]{chāzi — fork}

\label{lesson:ZH-C170-chazi}



\begin{quote}
Before the new one: say the Mandarin for a teacup, then the Mandarin for a spoon.
\end{quote}

\subsection*{You'll want to know: chāzi}
\emph{chāzi} — \textquotedblleft{}fork\textquotedblright{}.

\begin{quote}
\emph{nǐ yòng kuàizi háishi yòng chāzi} — Chopsticks or a fork?
\end{quote}

Its characters come later; for now it is a word for the ear.

\begin{cousinweb}
\emph{chāzi} is built from \emph{chā} and the noun ending \zh{子}, both inherited from Middle Chinese, \emph{tsrhae} and \emph{tsiX} (Guangyun, in Baxter's transcription). The dictionary glosses the first 'hands crossed', and under a second reading 'two branches': a forked thing. In the word the second loses its tone.
\end{cousinweb}

\begin{grammarlens}[title={What you've built}]
One more word for everyday things.
\end{grammarlens}

\subsection*{Your turn}
\begin{itemize}
\raggedright
  \item \emph{Say it:} \emph{chāzi}
  \item \emph{Say it:} \emph{chāzi}, once more
  \item \emph{Say it:} \emph{nǐ yòng kuàizi háishi yòng chāzi} — Chopsticks or a fork?
  \item \emph{Recall:} say the Mandarin for a teacup, then the Mandarin for a spoon, then say \emph{chāzi} again
\end{itemize}

\begin{tcolorbox}[breakable,colback=teal!4,colframe=teal!35!black,title={Before you move on}]
What does chāzi mean? (\textquotedblleft{}Fork\textquotedblright{}.) Say it once more.
\end{tcolorbox}
\section[wēibōlú]{wēibōlú — microwave oven}

\label{lesson:ZH-C170-weibolu}



\begin{quote}
Before the new one: say the Mandarin for a spoon, then the Mandarin for a fork.
\end{quote}

\subsection*{You'll want to know: wēibōlú}
\emph{wēibōlú} — \textquotedblleft{}microwave oven\textquotedblright{}.

\begin{quote}
\emph{fàn zài wēibōlú lǐ} — The food is in the microwave.
\end{quote}

Its characters come later; for now it is a word for the ear.

\begin{cousinweb}
\emph{wēibōlú} is built from \emph{wēi} 'tiny', \emph{bō} 'wave' and \emph{lú} 'stove', all three inherited from Middle Chinese, \emph{mj+j}, \emph{pa} and \emph{lu} (Guangyun, in Baxter's transcription): old words put together to name a new machine.
\end{cousinweb}

\begin{grammarlens}[title={What you've built}]
One more word for everyday things.
\end{grammarlens}

\subsection*{Your turn}
\begin{itemize}
\raggedright
  \item \emph{Say it:} \emph{wēibōlú}
  \item \emph{Say it:} \emph{wēibōlú}, once more
  \item \emph{Say it:} \emph{fàn zài wēibōlú lǐ} — The food is in the microwave.
  \item \emph{Recall:} say the Mandarin for a spoon, then the Mandarin for a fork, then say \emph{wēibōlú} again
\end{itemize}

\begin{tcolorbox}[breakable,colback=teal!4,colframe=teal!35!black,title={Before you move on}]
What does wēibōlú mean? (\textquotedblleft{}Microwave oven\textquotedblright{}.) Say it once more.
\end{tcolorbox}
